\subsection{Locked holdout}
\label{sec:results:holdout}

The supplemental window, 2021-12-06 to 2022-06-24 (\num{135} trading days),
was evaluated once per condition after that condition's design was frozen,
with each condition trained on all data to 2021-08-06 under the same validation
and embargo rules. It was opened for the eight network conditions, then for
the LightGBM baseline, and the eight network conditions were then re-run on it
under the released code because the first run could not be reproduced
(Section~\ref{sec:protocol:holdout}). Table~\ref{tab:holdout} reports the
reproducible run; Appendix~\ref{sec:app:rerun} gives the first.

\begin{table}[t]
\centering
\caption{Locked holdout, 135 trading days. Uncorrected
bootstrap $p$ against zero; Romano--Wolf $p$ across the eight contrasts with
the no-circuit control on this window (the two circuit-only
conditions and the second ansatz family were not); DSR deflated by the
\num{55} configurations (eleven models, five seeds) the study evaluated, using the
empirical dispersion of their Sharpe ratios on this window (sd \num{0.92},
expected maximum under the null \num{2.12}); break-even computed on the
pooled track, undefined where the gross return is not positive.}
\label{tab:holdout}
\small
\begin{tabular}{lrrrrrr}
\toprule
Condition & Sharpe & Seed sd & $p$ & $p_{\mathrm{RW}}$ & DSR & Break-even (bps) \\
\midrule
QTCNN, no circuit          & $+1.83$ & 0.84 & 0.07 & --   & 0.42 & 3.7 \\
QTCNN, classical surrogate & $+1.57$ & 1.08 & 0.18 & 0.95 & 0.34 & 5.5 \\
QTCNN, frozen circuit      & $+1.51$ & 1.05 & 0.17 & 0.95 & 0.33 & 4.4 \\
QTCNN, trained circuit     & $+1.49$ & 1.05 & 0.18 & 0.96 & 0.32 & 4.4 \\
TCN baseline               & $+1.37$ & 0.77 & 0.22 & 0.95 & 0.29 & 3.0 \\
Transformer baseline       & $+0.68$ & 0.87 & 0.47 & 0.54 & 0.15 & 2.1 \\
LSTM baseline              & $+0.67$ & 0.79 & 0.48 & 0.91 & 0.15 & 3.8 \\
LightGBM baseline          & $+0.49$ & 0.76 & 0.72 & 0.33 & 0.12 & 0.6 \\
MLP baseline               & $+0.42$ & 0.77 & 0.68 & 0.37 & 0.11 & 1.1 \\
\bottomrule
\end{tabular}
\end{table}

The window replicates the ablation and not the sign. Trained and frozen
circuits differ by \num{0.02}; the no-circuit control is the best of the four
hybrid conditions and the surrogate is no worse than either circuit; every
Diebold--Mariano contrast between a hybrid condition and the no-circuit control
has $p > 0.70$ (the classical baselines are closer, at \num{0.14} for LightGBM
and \num{0.16} for the MLP; one paired Sharpe-difference test, the MLP's,
reaches $p = 0.046$ uncorrected; none survives correction). Gross Sharpe
ratios are higher for all nine conditions than on the walk-forward track, which is a property of the
period, a stretch on which the cross-section rewarded the momentum and
liquidity features every model shares. The ranking of the nine is not stable
across the two windows: the transformer moves from first to sixth and the
no-circuit control from second to first, while the seed standard deviations
(\numrange{0.76}{1.08}) are wider than before because the window is short. What
is stable is the ablation itself (trained beside frozen, the surrogate no
worse than either, and no contrast significant), which is the quantity the
design was built to measure.

No condition is distinguishable from zero by the test on this window either.
For the best, the no-circuit control at \num{+1.83}, the percentile interval
just excludes zero ($[+0.02, +3.94]$) while the null-imposed bootstrap gives
$p = 0.07$, and its Deflated Sharpe Ratio is \num{0.42} against the
\num{0.95} a claim would require: deflated by the \num{55} configurations, the
expected maximum Sharpe under the null on \num{135} days is \num{2.12}, above
every row of the table. At \SI{5}{bps} one condition remains positive, the
surrogate at \num{+0.18}; at \SI{10}{bps} none do.
